\section{System overview}
\label{sec:visual_tour}
The project simulates an automated warehouse sorting system using OpenGL. Boxes enter on an upper conveyor, fall onto a lower conveyor, stop for a simulated scan, and follow one of two exits. Blue boxes are directed to bin A, while orange boxes are directed to bin B.
\begin{shot}{assets/overview.png}{The complete rendered factory. Labels identify the sorting machine and the surrounding objects.}
\callout{0.47,0.56}{0.19,0.80}{Sorting machine}
\callout{0.35,0.48}{0.18,0.34}{Parked forklift}
\callout{0.67,0.59}{0.78,0.82}{Storage racks}
\callout{0.59,0.445}{0.83,0.30}{Control cabinet}
\callout{0.69,0.69}{0.87,0.66}{Wall fan}
\end{shot}
\subsection{Sorting sequence}
\begin{enumerate}
\item \textbf{Entry and transport:} colored cuboid parcels travel on the upper conveyor.
\item \textbf{Transfer and scanning:} a parcel falls onto the lower belt, then waits 0.65 s at the scanner.
\item \textbf{Routing and collection:} the gate aligns with the assigned route. The parcel turns, drops into its bin, and increments the corresponding counter.
\end{enumerate}
\note{The scanner is simulated, and the destination is assigned when a parcel is spawned. The gate and parcel trajectory use coordinated animations.}
\subsection{Technical specifications}
\par{\small\begin{tabularx}{\textwidth}{@{}L{34mm}Y@{}}\toprule
Software & C++17, OpenGL 3.3 Core requested, GLSL 330, CMake 3.20+.\\
Dependencies & GLFW 3.3.9, GLAD 2.0.6, GLM 0.9.9.8.\\
Scene construction & Reusable cubes and cylinders; no imported factory model or image texture.\\
Simulation and rendering & 12 reusable parcel slots, four point lights, daylight and shadows.\\\bottomrule
\end{tabularx}}\par
